Verteilte Satellitensysteme k"onnen zeitkritische Erdbeobachtungsaufgaben schneller bearbeiten, ihre Leistungsf"ahigkeit h"angt jedoch von orbitalem Zugang, begrenzten Kommunikationsfenstern und der gew"ahlten Weiterleitungsstrategie ab. Diese Arbeit untersucht den Einfluss eines Zentralenknotenanteils auf die Verteilung von Waldbrandaufgaben, die nachfolgende Beobachtung, den Ressourcenbedarf und die Empfindlichkeit gegen Ausf"alle. Eine nachvollziehbare Datenpipeline erzeugt aus deduplizierten Sentinel-3-SLSTR-FRP-Detektionen 805 priorisierte Aufgaben: 32 f"ur Kalifornien und 773 f"ur Indien.

Der Entwurfsraum umfasst 60, 120 und 180 Satelliten, 3, 5 und 10 Bahnebenen, H"ohen von 500, 800 und 1000 km, Inklinationen von 60, 80 und 98,6 Grad sowie Zentralenknotenanteile zwischen 0 und 100 Prozent. Die fr"uhere nominale Kampagne enth"alt 891 validierte F"alle pro Region. Die neuere Full891-Fresh-V2-Kampagne definiert 539 Betriebs-, Policy-, Aufgabengr"o\ss en- und Fehlerszenarien pro Architektur. F"ur Kalifornien liegen 890 vollst"andige F"alle und nach Entfernung von Resume-Duplikaten 480.110 eindeutige Szenariozeilen vor; alle 19.179 Architekturpr"ufungen wurden bestanden.

In der begrenzten nominalen Policy betr"agt die mittlere Beobachtungserf"ullung 61,59 Prozent, die strikte vollst"andige Verteilung an n"utzliche Empf"anger 0,63 Prozent und deren mittlere Abdeckung 15,29 Prozent. Das diagnostische Szenario ohne Hop- und Kopienbegrenzung erh"oht diese Werte auf 77,62, 74,69 und 91,14 Prozent, steigert jedoch die mittlere "ubertragene Datenmenge von 197,55 auf 2.047,57 Mbit pro Fall. Die Wirkung des Zentralenknotenanteils ist nicht monoton und zeigt abnehmenden Grenznutzen bei weiter wachsender Kommunikationslast.

Bei einem Stressniveau von 30 Prozent ist der Ausfall der einzigen Bodenstation die dominante kalifornische Schwachstelle. Eine Vervierfachung der Aufgabengr"o\ss e reduziert die mittlere Beobachtungserf"ullung von 77,62 auf 69,48 Prozent. Eine Pareto-konsistente Gewichtungsanalyse bevorzugt bei ausgewogenen Annahmen eine Konfiguration mit 60 Satelliten, 10 Ebenen, 800 km H"ohe, 98,6 Grad Inklination und 40 Prozent Zentralenknoten. Diese Konfiguration gewinnt 41,79 Prozent von 10.000 zuf"alligen Gewichtungsziehungen. Das Ergebnis ist ein bedingter Kompromiss und kein universelles Optimum. Die Indien-Fresh-V2-Ergebnisse und eine Sensitivit"at des Beobachtungsradius bleiben vor der abschlie\ss enden regions"ubergreifenden Robustheitsaussage offen.
